\documentclass[12pt]{article}

\usepackage[utf8]{inputenc}
\usepackage{amsmath, amssymb, amsthm}
\usepackage{geometry}
\geometry{a4paper, margin=2.5cm}

\title{Teorema de Structură a Modulelor Finit Generate peste un PID}
\author{}
\date{}

\theoremstyle{definition}
\newtheorem{definition}{Definiție}

\theoremstyle{plain}
\newtheorem{theorem}{Teoremă}

\begin{document}

\maketitle

\begin{definition}
Pentru un $R$-modul $M$ și un element $x \in M$, \emph{anihilatorul} lui $x$ este idealul


\[
\operatorname{ann}(x) = \{\, r \in R : r x = 0 \,\}.
\]


\end{definition}

\begin{theorem}[Teorema de structură pentru module finit generate peste un PID]
Fie $R$ un domeniu principal (PID) și $M$ un $R$-modul finit generat. Atunci există
elemente $x_1, \dots, x_m \in M$ și elemente nenule $d_1, \dots, d_m \in R$ cu


\[
d_1 \mid d_2 \mid \cdots \mid d_m
\]


astfel încât


\[
M \cong R/d_1 R \;\oplus\; R/d_2 R \;\oplus\; \cdots \;\oplus\; R/d_m R \;\oplus\; R^k,
\]


unde $k$ este rangul liber al lui $M$.

Mai mult, putem alege $x_i$ astfel încât


\[
\operatorname{ann}(x_i) = d_i R.
\]


\end{theorem}

\begin{theorem}[Forma în factori invarianti]
În ipotezele teoremei, există o bază a unui modul liber $F$ și un submodul $E$ astfel încât


\[
M \cong F/E,
\]


iar $E$ este generat de elementele


\[
d_1 u_1,\; d_2 u_2,\; \dots,\; d_m u_m,
\]


cu $d_1 \mid d_2 \mid \cdots \mid d_m$.
\end{theorem}

\begin{theorem}[Forma în factori elementari]
Fie $d_i = p_{i1}^{e_{i1}} \cdots p_{ik_i}^{e_{ik_i}}$ factorizarea în factori primi.
Atunci


\[
M \cong \bigoplus_{i=1}^m \bigoplus_{j=1}^{k_i} R / p_{ij}^{\, e_{ij}} R.
\]


Aceasta este descompunerea în factori elementari.
\end{theorem}

\end{document}
